
\section{11. Conclusion}

\subsection{11.1 Summary}

We introduced \textbf{Primal Attention}, an asymmetric kernel SVD formulation 
of self-attention. By allowing different landmarks for queries and keys, 
Primal Attention generalizes symmetric Nystr m approximation. The pseudoinverse 
is computed via differentiable Schulz iteration.

\subsection{11.2 Results Recap}

\begin{table}[H]
\centering
\small
\begin{tabular}{lll}
\toprule
\textbf{Model} & \textbf{Acc $\Delta$} & \textbf{Attn mem ratio} \\
\midrule
DistilBERT & -4.2\% & **4$\times$** \\
RoBERTa & -4.6\% & **4$\times$** \\
GPT-2 & -4.4\% & **4$\times$** \\
\bottomrule
\end{tabular}
\end{table}

\subsection{11.3 Key Takeaways}

1. \textbf{Asymmetric landmarks} give +0.5\% over symmetric Nystr m.
2. \textbf{$m = 64$} is optimal for $L = 128$.
3. \textbf{Schulz from SVD pinv} is best practice.
4. \textbf{Long $L$ benefits most.}
5. \textbf{No parameter reduction.}

\subsection{11.4 Future Work}

\subsubsection{Short-Term}
\begin{itemize}
\item True causal approximation
\item Learned landmark pooling
\item Fused kernels
\item Long-context evaluation
\end{itemize}

\subsubsection{Long-Term}
\begin{itemize}
\item Theoretical error bounds
\item Integration with MoE
\item Cross-task transfer
\end{itemize}

\subsection{11.5 Closing Remarks}

Primal Attention provides a flexible approximation primitive for attention. 
We release code and checkpoints.

\subsection{11.6 Acknowledgments}

We thank {{ACKNOWLEDGMENTS}}.

\subsection{11.7 Author Contributions}

\begin{itemize}
\item \textbf{{{AUTHOR\_1}}} designed Primal Attention, implemented the code, wrote the paper.
\item \textbf{{{AUTHOR\_2}}} provided feedback, edited.
\end{itemize}

\subsection{11.8 Conflict of Interest}

None.